% Generated draft: requires mathematical and source review.
\documentclass[11pt]{article}
\usepackage[margin=25mm]{geometry}
\usepackage{fontspec}
\setmainfont{Latin Modern Roman}
\usepackage{amsmath,amssymb,mathtools,unicode-math}
\setmathfont{Latin Modern Math}
\usepackage{enumitem,needspace,xurl,hyperref}
\setlength{\parindent}{0pt}
\setlength{\parskip}{5pt}
\setlength{\emergencystretch}{3em}
\newcommand{\N}{\mathbb{N}}
\newcommand{\Z}{\mathbb{Z}}
\newcommand{\Q}{\mathbb{Q}}
\newcommand{\R}{\mathbb{R}}
\newcommand{\C}{\mathbb{C}}
\newcommand{\F}{\mathbb{F}}
\newcommand{\A}{\mathbb{A}}
\newcommand{\PP}{\mathbb{P}}
\newcommand{\E}{\mathbb{E}}
\newcommand{\eps}{\varepsilon}
\newcommand{\dd}{\,\mathrm{d}}
\begin{document}
\section*{UnsolvedMath ID: 3100039}
\textbf{Source problem code:} catalogue entry
\subsection{Definitions and mathematical statement}
Let $S_n$ be the permutations of $\{1,\ldots,n\}$, and let
\[
 \ell(w)=\#\{(i,j):1\leq i<j\leq n,\ w(i)>w(j)\}
\]
be inversion number. Fix the right weak order: its cover relations are $w\lessdot ws_i$ when $s_i=(i,i+1)$ is an adjacent transposition and $\ell(ws_i)=\ell(w)+1$. The order is the transitive closure of these covers, ranked by $\ell$. The left weak order is isomorphic by inversion, so this choice does not change the Sperner question. No strong Bruhat order is substituted.

For $0\leq r\leq\binom n2$, put $L_r=\{w\in S_n:\ell(w)=r\}$. An antichain is a subset with no two distinct comparable permutations. A ranked poset is Sperner when its maximum antichain size equals its largest rank size.

\textbf{Question.} For every $n\geq1$, is weak order on $S_n$ Sperner; equivalently, does every antichain $A$ satisfy
\[
 |A|\leq\max_{0\leq r\leq\binom n2}|L_r|?
\]
The reverse inequality for the maximum is represented by any largest rank, since each rank is an antichain.

The supplied summary marks this solved and further claims strong Sperner and normalized-matching properties. Strong Sperner would control a union of any $k$ antichains by the sum of the $k$ largest rank sizes, a stronger assertion than the recorded question. The summary attributes these claims to Stanley, ``Weyl groups, the hard Lefschetz theorem and the Sperner property'' (1980), but gives no theorem number checking that the cited order is the weak order. Its tentative parenthetical about ``maximum-length elements'' does not define largest rank: rank $\binom n2$ consists of the unique longest permutation, whereas largest rank means largest cardinality among all $L_r$. The cited order, stronger properties, and attribution require verification; none is independently certified here.
\subsection{Short English statement}
Is every antichain in permutation weak order no larger than the largest inversion-number level? The record says yes and claims stronger properties, but its Stanley citation and tentative rank description need checking, especially the distinction between weak and strong Bruhat orders.
\subsection{Sources}
\begin{itemize}
\item Local UnsolvedMath registry, permanent ID 3100039, source code catalogue entry, ``Is the weak order on S\_(n) (the inversion poset) Sperner''; status solved, classification SOLVED-IN-LITERATURE, updated 2026-08-06.
\item Cooper, \emph{Combinatorial Problems I Like}, October 2020 snapshot, section ``Matroids, Lattices, and Posets,'' source-order bullet 39, \url{https://people.math.sc.edu/cooper/combprob.html}. The record reports HTML extraction and access on 2026-07-29. It says later resolution was not checked item by item. Rights: NEEDS\_REVIEW; public access verified but redistribution permission not established.
\item Stanley, ``Weyl groups, the hard Lefschetz theorem and the Sperner property'' (1980), as supplied in the research summary. No journal, pages, theorem number, or explicit weak-order identification is given; the strong-Sperner and normalized-matching claims are also summary claims.
\end{itemize}
\end{document}
